% TOP_PROBLEM: {"id": 2800203, "title": "10 Lectures and 42 Open Problems — The planted clique problem", "category_id": 3, "category": {"id": 3, "name": "graph_theory", "display_name": "Graph Theory", "description": "Problems involving graphs, networks, and their properties.", "slug": "graph-theory", "order_index": 3, "created_at": "2026-07-31T15:26:25.671Z"}, "status": "open", "problem_number": "AMR-027-0203", "set_id": 13}
% FIRST_POSTED: 2026-10-02
% ENTRY_KIND: statement_only
% UnsolvedMath ID 2800203; source catalogue code AMR-027-0203
% !TeX program = lualatex
% Definition and sources only; this file is not an LLM solution attempt.
\documentclass[11pt,a4paper]{article}
\usepackage[margin=25mm]{geometry}
\usepackage{fontspec}
\setmainfont{lmroman10-regular.otf}[BoldFont=lmroman10-bold.otf,
 ItalicFont=lmroman10-italic.otf,BoldItalicFont=lmroman10-bolditalic.otf]
\usepackage{amsmath,amssymb,mathtools}
\usepackage{unicode-math}
\setmathfont{latinmodern-math.otf}
\AtBeginDocument{\let\setminus\smallsetminus}
\usepackage{xurl,hyperref}
\hypersetup{colorlinks=true,urlcolor=blue}
\setcounter{secnumdepth}{0}
\setlength{\parindent}{0pt}
\setlength{\parskip}{5pt}
\setlength{\emergencystretch}{3em}
\begin{document}
\section*{10 Lectures and 42 Open Problems — The planted clique problem}\label{2800203}
{\small UnsolvedMath ID 2800203 · AMR-027-0203 · Graph Theory · AMR Open Problem Lists}
\subsection{Definitions and mathematical statement}
On the labelled vertex set $[n]=\{1,\ldots,n\}$, first include each
unordered pair independently with probability $1/2$. Independently choose
a uniformly random set $S\subseteq[n]$ of size $k$ and add all missing
edges within $S$. Denote the resulting distribution by $\mathcal G_{n,k}$;
the null distribution $\mathcal G_{n,0}$ has no planted set.

The recovery problem asks an algorithm to output the planted clique,
or equivalently the maximum clique in the source's regime
$k\gg\log n$. The detection problem asks it to distinguish these two
graph distributions. Success with high probability means probability
$1-o(1)$ as $n\to\infty$, including graph and algorithm randomness;
for detection, this requirement applies under both hypotheses.

Can either task be solved in time polynomial in $n$ when
\[
                 k=\left\lfloor\frac{\sqrt n}{\log n}\right\rfloor,
\]
or more generally in a nontrivial regime $\log n\ll k\ll\sqrt n$?
The polynomial exponent must be independent of $n$.

The source also asks whether recovery is possible in quasi-linear time
in the dense input size, namely $n^2(\log n)^{O(1)}$, below the threshold
$k=(1/\sqrt e-\varepsilon)\sqrt n$ for some fixed
$0<\varepsilon<1/\sqrt e$. Finding a large clique in an arbitrary graph
is a different worst-case problem. Here the entire promise is the
specified random construction; the algorithm is not given $S$.
\subsection{Short English statement}
Can planted cliques below the square-root scale be detected or recovered in polynomial time, or recovered below the stated threshold in quasi-linear dense-input time?
\subsection{Sources}
\begin{itemize}
\item[S1] ULAM.AI and the UnsolvedMath Contributors, supplied problems.json, record 2800203 (AMR-027-0203), AMR Open Problem Lists. Catalogue identity and source transcription; CC BY 4.0.
\url{https://huggingface.co/datasets/ulamai/UnsolvedMath}.
\item[S2] Bandeira - 42 Open Problems in Data Science (2016), Open Problem 2.3. Original source indexed by AMR-027-0203.
\url{https://afonsobandeira.wordpress.com/2015/09/30/10l42pgraphs/}.
\item[S3] A. S. Bandeira, Ten Lectures and Forty-Two Open Problems in the Mathematics of Data Science, corresponding numbered problem and surrounding definitions.
\url{https://people.math.ethz.ch/~abandeira/TenLecturesFortyTwoProblems.pdf}.
\end{itemize}
\end{document}
